\documentclass[10pt,a4paper]{amsart}
\usepackage[margin=19mm]{geometry}
\usepackage{amsmath,amssymb,mathtools}
\usepackage[scr=rsfs,cal=euler]{mathalfa}
\usepackage{xfrac}
\usepackage[unicode,hidelinks,bookmarks=false]{hyperref}
\newcommand{\ie}{i.e.\ }
\newcommand{\cf}{cf.\ }
\newcommand{\resp}{respectively\ }
\newcommand{\Subsection}{\subsection}
\providecommand{\nobreakdash}{\nobreak\hbox{-}}
\setlength{\emergencystretch}{2em}
\multlinegap0pt
\usepackage{amssymb}
\usepackage{enumerate}
\newtheorem{lemma}{Lemma}
\DeclareMathOperator{\Ca}{Cap}
\renewcommand{\P}{\mathbb{P}}
\newcommand{\Z}{\mathbb{Z}}
\title{Capacity of loop-erased random walk}
\author{Maarten Markering}
\date{Source: arXiv:2411.13505v3 (CC BY 4.0)}
\begin{document}
\maketitle

\noindent\emph{Source and licence.} Capacity of loop-erased random walk. Version arXiv:2411.13505v3 (CC BY 4.0), distributed by arXiv under the Creative Commons Attribution 4.0 International licence, http://creativecommons.org/licenses/by/4.0/, as recorded in the arXiv licence metadata for this exact version and shown on the abstract page https://arxiv.org/abs/2411.13505v3. Full licence notice: \texttt{licenses/arxiv-2411.13505-CC-BY-4.0.txt}.\par\smallskip

\noindent\textit{Editorial scope.} Counted unit: the article's lemma lem:capacitydecomp (source0.tex lines 462--467). The article's complete original proof of it is delivered with it (source0.tex lines 468--488, one \texttt{proof} environment of 21 lines). No mathematics is written, repaired or re-derived: the statement and the proof are the article's own lines, in the article's own notation, and no retained line is substituted or re-flowed.  No further source-local context is needed: every cross-reference of every retained line resolves inside the two delivered spans.  Nothing was omitted from any retained span, so the package carries no omission record.  No retained line contains a citation, so the wrapper carries no bibliography and no bibliography entry.  No retained line contains a figure environment, an image-inclusion command or drawing code, so no image is delivered and the package declares no asset dependency.  Editorial formatting only, in addition: an AMS \texttt{amsart} preamble that restates the article's own declarations and macros used by the retained lines, namely the environments \texttt{lemma} on separate counters, together with the standard packages those lines need.  The retained environments print in this document as Lemma 1 in order of appearance, whereas the article prints the same items with the numbering of its own full document.  The complete native source of the article as distributed in the arXiv e-print for this exact version is bundled unchanged as \texttt{sources/168826/source0.tex}.
\begin{lemma}\label{lem:capacitydecomp}
Let $A=\{x_1,\ldots,x_n\}\subset\Z^d$. Then
\begin{equation}
    \Ca(A)=\sum_{k=1}^n\P_{x_k}(W[1,\infty)\cap\{x_1,\ldots,x_k\}=\varnothing)\P_{x_k}(W[1,\infty)\cap\{x_1,\ldots,x_{k-1}\}=\varnothing).
\end{equation}
\end{lemma}

\begin{proof}
For $B\subset\Z^d$, let $\tau_B:=\min\{t\geq0\colon W(t)\in B\}$ be the first hitting time of $B$. The proof is by induction. The statement is clearly true for all sets of size 1. Now assume the statement holds for all sets of size $n-1$. Then
\begin{equation}
    \begin{split}
        \Ca(A)=&\sum_{k=1}^n\P_{x_k}(W[1,\infty)\cap\{x_1,\ldots,x_n\}=\varnothing)\\
        =&\sum_{k=1}^n[\P_{x_k}(W[1,\infty)\cap\{x_1,\ldots,x_{n-1}\}
        =\varnothing)\\
        &-\P_{x_k}(W[1,\infty)\cap\{x_1,\ldots,x_{n-1}\}=\varnothing,\,W[1,\infty)\cap\{x_n\}\neq\varnothing)]\\
        =&\Ca(\{x_1,\ldots,x_{n-1}\})+\P_{x_n}(W[1,\infty)\cap\{x_1,\ldots,x_{n-1}\}=\varnothing)\\
        &-\sum_{k=1}^n\P_{x_k}(\tau^+_{x_n}=\tau^+_{\{x_1,\ldots,x_{n}\}}<\infty)\P_{x_n}(W[1,\infty)\cap\{x_1,\ldots,x_{n-1}\}=\varnothing)\\
        =&\Ca(\{x_1,\ldots,x_{n-1}\})\\
        &+\P_{x_n}(W[1,\infty)\cap\{x_1,\ldots,x_{n-1}\}=\varnothing)\left(1-\sum_{k=1}^n\P_{x_n}(\tau_{x_k}^+=\tau_{\{x_1,\ldots,x_n\}}^+<\infty)\right)\\
        =&\Ca(\{x_1,\ldots,x_{n-1}\})\\
        &+\P_{x_n}(W[1,\infty)\cap\{x_1,\ldots,x_{n-1}\}=\varnothing)(1-\P_{x_n}(W[1,\infty)\cap\{x_1,\ldots,x_{n}\}\neq\varnothing))\\
        =&\Ca(\{x_1,\ldots,x_{n-1}\})\\
        &+\P_{x_n}(W[1,\infty)\cap\{x_1,\ldots,x_{n-1}\}=\varnothing)\P_{x_n}(W[1,\infty)\cap\{x_1,\ldots,x_{n}\}=\varnothing)\\
        =&\sum_{k=1}^n\P_{x_k}(W[1,\infty)\cap\{x_1,\ldots,x_k\}=\varnothing)\P_{x_k}(W[1,\infty)\cap\{x_1,\ldots,x_{k-1}\}=\varnothing),
    \end{split}
\end{equation}
which completes the proof.
\end{proof}

\end{document}
